Consistent with this downstream focus, the error profiles placed the two baselines at opposite operating-point extremes, whereas the proposed methods provided a closer balance between false positives and false negatives. BLINKER-concat was strongly false-positive dominated, with an FP:FN ratio of 26.675, whereas MNE-annot was false-negative dominated, with a ratio of 0.613. By contrast, Proposed-Mean and Proposed-Med produced ratios of 1.206 and 1.628, respectively, indicating substantially closer correspondence between over-detection and missed events. This balance matters because methods with similar $F_1$ can impose different downstream costs: excessive false positives inflate estimated blink counts and increase the burden of rejecting spurious regions, whereas excessive false negatives suppress genuine events and bias analyses that depend on blink occurrence or timing. An operating point closer to parity therefore provides a more usable compromise when neither error type can be treated as negligible. Precision--recall characterisation is the appropriate framework for this comparison because blink detection is an imbalanced event-detection problem in which the positive class is sparse relative to continuous background signal~\citep{saito2015precision}. Event-level evaluation is also preferable to sample-wise agreement because it assesses whether temporally localised detections correspond to distinct underlying events~\citep{salles2024softed}.
